% TCC — Introdução
\section{Introdução}
APIs REST tornaram-se parte essencial da infraestrutura de sistemas modernos, conectando serviços, aplicações web e dispositivos móveis \cite{arcuri2026fuzzingrestapisindustry}. Nesse cenário, falhas de robustez e inconsistências entre o comportamento da API e o contrato publicado impactam diretamente a confiabilidade de integrações e o custo de manutenção. Embora testes funcionais tradicionais sejam eficazes para validar cenários felizes (\textit{happy path}), eles frequentemente deixam de exercitar valores-limite, formatos inválidos e combinações pouco usuais de parâmetros, que são justamente as situações em que falhas sutis tendem a se manifestar.

O uso de contratos formais, como a especificação OpenAPI (versões 3.0 e 3.1), viabiliza uma abordagem "guiada por contrato" para a geração de casos de teste: em vez de produzir entradas aleatórias, o processo explora estruturas, tipos e restrições descritas no \textit{schema} para produzir mutações \emph{válidas e inválidas} de forma sistemática. Essa estratégia aumenta a probabilidade de revelar respostas inesperadas (por exemplo, códigos 5xx, mensagens de erro inconsistentes ou divergências estruturais entre a resposta e o contrato) sem exigir acesso ao código-fonte nem instrumentação interna da aplicação, isto é, em um contexto \emph{black-box} \cite{meyer1992designbycontract}.

Este trabalho propõe o desenvolvimento e a avaliação de uma ferramenta de \emph{fuzzing} automatizado para APIs REST baseada em contratos OpenAPI. A ferramenta implementa: (i) um \textit{parser} que converte um subconjunto da especificação em um modelo interno; (ii) um motor que produz entradas-base e aplica mutações orientadas por tipos e por parte das restrições representadas; (iii) um plano de execução e um \textit{runner} de requisições com controle de concorrência; e (iv) um analisador que registra evidências e identifica divergências em relação ao contrato interpretado. A avaliação caracteriza os achados produzidos em APIs distintas, sua reprodução entre campanhas e a utilidade do mecanismo de priorização.

Além da identificação de divergências contratuais, a ferramenta utiliza um sistema de inferência fuzzy do tipo Mamdani para atribuir um escore contínuo aos achados reportados. Esse componente não confirma automaticamente defeitos ou vulnerabilidades: sua função é apoiar a priorização da inspeção humana a partir de sinais observáveis da execução. Nesse contexto, a pesquisa investiga quais tipos de achado o \textit{fuzzer} produz em APIs distintas, em que medida esses achados são reproduzíveis e quão coerente é o ordenamento de severidade em relação a critérios humanos definidos independentemente do escore.

O escopo delimita-se a APIs REST descritas em OpenAPI, não abrangendo cenários dependentes de sequenciamento complexo de requisições ou manutenção de estado entre chamadas, autenticação avançada (por exemplo, fluxos completos de OAuth2) nem avaliação de carga ou desempenho. A ferramenta gera saídas estruturadas compatíveis com automação, mas a integração completa e a efetividade em fluxos de CI/CD não constituem objetos centrais da avaliação.

O restante deste artigo está organizado da seguinte forma. A Seção~\ref{sec:fundamentacao} apresenta os conceitos de OpenAPI, \textit{fuzzing} de APIs e lógica fuzzy necessários à compreensão do trabalho. A Seção~\ref{sec:relacionados} discute ferramentas e pesquisas relacionadas e posiciona a proposta. A Seção~\ref{sec:metodologia} descreve o delineamento e o protocolo de avaliação, enquanto a Seção~\ref{sec:desenvolvimento} detalha a arquitetura e as decisões da implementação. A Seção~\ref{sec:resultados} apresenta e discute os resultados experimentais. Por fim, a Seção~\ref{sec:conclusoes} sintetiza as contribuições, limitações e possibilidades de continuidade.
